\section{本章小结}
\label{sec:ctsp-summary}

本章开头的三台记录仪分别给出跳跃路径、连续抖动的路径，以及噪声与状态相互作用的路径。回头看，它们沿着“路径—条件未来—累计与停止—状态方程—分布与长期行为—近似”这条线得到了处理，每一步都与上一章的离散时间结论对应，又都补上了连续时间特有的条件。

路径层面，Poisson过程由独立指数等待时间构造，Brownian运动由独立高斯增量与连续版本定义。时间连续并不意味着路径连续：前者是右连续有左极限的跳跃路径，后者处处连续而总变差无限、二次变差等于时间。不可数时刻带来的问题由三项约定解决：取右连续有左极限的版本，使修改与不可区分在其中一致；使用满足通常条件的滤过；以每个实数时刻上可判断的停时，并用二进上逼近把离散结论搬过来。

条件结构有两种，它们可以在同一个过程上同时成立。Markov性质由转移半群描述，半群又由生成元这一短时间变化率决定；连续时间Markov链把这些对象变成矩阵，但可数状态时还要检查不爆炸，强Markov性质还需要路径与Feller正则性。鞅性质约束条件均值；可选抽样、Doob最大不等式与鞅收敛都在右连续版本上成立，无界停时仍要一致可积，局部鞅取期望之前要验证可积性。计数过程扣除可预测的补偿量后才是鞅。

累计与状态方程是本章新增的核心。对时间与有限次跳跃可以逐路径累计，跳跃系数必须可预测；对Brownian运动只能借Itô等距在$L^2$中构造积分。二次变差使Itô公式多出$\frac12\operatorname{tr}(\bm\sigma\bm\sigma^\top\nabla^2V)$，跳跃则保留完整的有限差；Stratonovich约定恢复普通链式法则，但与Itô约定之间的转换要修正漂移。随机微分方程在全局Lipschitz条件下由Picard迭代得到强解，局部条件下需要Lyapunov函数排除爆炸；几何Brownian运动、OU过程与线性系统给出了可以直接核对的显式解。

从路径到整体规律的桥梁是生成元。Dynkin公式与含时间的Itô公式把期望、密度与首次到达化为后向方程、Fokker--Planck方程与边值问题，Feynman--Kac公式又把一类线性PDE的解写成路径期望。长期行为需要区分两个问题：同一噪声下的路径是否稳定，以及分布是否趋于不变分布。乘性噪声说明几乎必然稳定与均方稳定可以不一致；形式稳态密度只是候选，存在、唯一与收敛要分别论证；时间平均的误差由长期方差或Poisson方程控制，遍历性本身不保证中心极限定理。

表~\ref{tab:ctsp-discrete-continuous}把这些对应关系集中列出。

\begin{table*}[htbp]
  \centering
  \begin{tabularx}{\linewidth}{@{}l>{\RaggedRight\arraybackslash}p{50mm}>{\RaggedRight\arraybackslash}X@{}}
    \toprule
    问题 & 离散时间（第~\ref{chap:discrete-time-stochastic-processes}章） & 连续时间新增的对象与条件\\
    \midrule
    怎样生成路径 & 转移矩阵、转移核、随机递推 & 等待时间与跳转率；Brownian增量与连续版本；随机微分方程\\
    条件未来 & 多步核、Chapman--Kolmogorov & 转移半群与生成元的定义域；不爆炸；强Markov需路径与Feller正则\\
    累计与停止 & 鞅变换、可选抽样、Doob不等式 & 右连续版本、二进逼近、局部化；可预测补偿；Itô积分与二次变差\\
    分布演化 & 分布右乘转移核 & 弱形式、Fokker--Planck方程（需密度正则）、后向方程与边值问题\\
    长期行为 & 不变分布、遍历平均、相关误差 & 半群不动点、骨架链与连续过程的区别、Poisson方程与长期方差\\
    近似与计算 & 有限步偏差、平均误差 & 强误差与弱误差、事件驱动模拟、步长偏差、路径空间紧性\\
    \bottomrule
  \end{tabularx}
  \caption{离散时间与连续时间随机过程的对应。右列只列出连续时间新增的对象与条件，离散时间的定义与结论在相应位置简短回顾后沿用。}
  \label{tab:ctsp-discrete-continuous}
\end{table*}

最后三类计算任务使用相同的工具，目标却不同，应当分开理解。数值模拟忠实地近似一个给定的模型，误差按路径或期望衡量，并且有限时段的误差界不能直接推出长期行为。采样积分设计一个以目标分布为不变分布的过程，要求它长期保持非退化的随机性，按停留时间累计平均；离散化会改变不变分布，需要接受校正或步长控制。连续时间随机逼近则让增益递减、噪声逐渐压低，使状态收敛到指定的根；在算法时钟下它接近平均ODE，放大波动后又呈现扩散的形态。Itô积分是这三者共同的分析工具，而不是它们要计算的数值积分。

路径极限理论为离散近似与连续模型之间的关系提供了最后的依据：有限维分布收敛加上路径紧性，才能推出整条路径的分布收敛，并经连续映射传给最大值、时间积分与首次到达等泛函。
